\documentclass[11pt]{article}
\usepackage[margin=2.5cm]{geometry}
\usepackage{amsmath,amssymb,amsthm}

\newtheorem{lemma}{Lemma}
\renewcommand{\thelemma}{R\arabic{lemma}}
\theoremstyle{remark}
\newtheorem*{remark}{Remark}

\newcommand{\N}{\mathbb{N}}
\newcommand{\diag}{\operatorname{diag}}

\title{Problem B, Cell 3 (A General Upper Bound): partial submission\\[2pt]
\large Lower bound $D_B(T_k-1)\ge k^2-2k-1$ for every $k\ge 3$}
\author{Team submission (Scribe B-C3-014), from gated artefact B-C3-002}
\date{}

\begin{document}
\maketitle

\section*{Statement (verbatim)}

\textbf{Definitions (verbatim, problem statement).} A \emph{partition} of a positive integer $n$ is a weakly
decreasing sequence $\lambda=(\lambda_1,\ldots,\lambda_s)$ of positive integers with $\lambda_1+\cdots+\lambda_s=n$.
Think of it as a division of $n$ cards into $s$ piles.
The \emph{shift} $B(\lambda)$ is the partition of $n$ obtained as follows: remove one card from every pile,
discard the piles that become empty, and add one new pile consisting of the $s$ removed cards. Formally,
$B(\lambda)$ is the partition whose parts are the positive numbers among $\lambda_1-1,\ldots,\lambda_s-1$ together
with one extra part equal to $s$.
Call $\lambda$ \emph{cyclic} if $B^i(\lambda)=\lambda$ for some $i\ge 1$, and let
\[
d_B(\lambda)=\min\{\,i\ge 0: B^i(\lambda)\text{ is cyclic}\,\},\qquad
D_B(n)=\max\{\,d_B(\lambda):\lambda\text{ is a partition of } n\,\}.
\]
Write $T_k=\frac{k(k+1)}{2}$, $\delta_k=(k,k-1,\ldots,2,1)$. Every $n\ge 1$ satisfies $T_{k-1}<n\le T_k$ for
exactly one $k$, which we call the \emph{rank} of $n$.

\medskip
\textbf{C3: A General Upper Bound (verbatim).} Now the numbers strictly between two consecutive triangular numbers.
\begin{itemize}
\item[(a)] Prove that for every $k\ge4$ and every non-triangular $n$ with $T_{k-1}<n<T_k$,
\[ D_B(n)\le k^2-2k-1. \]
\item[(b)] Determine $D_B(T_k-1)$ exactly.
\item[(c)] Determine also, for that $n$, which partitions attain the maximum.
\end{itemize}

\section*{Status}

\begin{itemize}
\item \textbf{Cell status:} PARTIAL.
\item \textbf{Claim status of the established claim:} PROVED.
\item \textbf{Established.} (b) LOWER half: for every $k\ge 3$,
$\lambda^*_k=(k-1,k-2,k-2,k-3,\ldots,2,1,1)$ is a partition of $T_k-1$ with $d_B(\lambda^*_k)=k^2-2k-1$, hence
$D_B(T_k-1)\ge k^2-2k-1$ (Sections~0--6 below; PROVED).
\item \textbf{Remaining gap (OPEN).} (a) the upper bound $d_B\le k^2-2k-1$ for every non-triangular $n$ of rank
$k\ge4$; the upper half of (b) (so $D_B(T_k-1)=k^2-2k-1$ is \emph{not} established); (c) the maximiser set.
All three are OPEN. See the section ``Open parts''.
\end{itemize}

\section*{Cited versus own work}

\begin{itemize}
\item \textbf{Cited.} The Cell 1 classification of cyclic partitions (the earlier cell C1 of this problem, gated
earlier in this run; it is used as a cited result and is not re-proved here). In the form used below: for $n$ of
rank $k$, $n=T_{k-1}+r$ with $1\le r\le k$, the cyclic partitions of $n$ are exactly
\[
\lambda(e)=(k-1+e_1,\ k-2+e_2,\ \ldots,\ 1+e_{k-1},\ e_k),\qquad e\in\{0,1\}^k,\ \textstyle\sum e=r
\]
(zero entries omitted). It is used only in Lemma~R2 (Section~2.5).
\item \textbf{Own work (team).} Sections~0--6 (Lemmas R1--R6): the rotation lemma, the energy characterisation
of cyclicity, the classification and exact dynamics of partitions with energy excess $\Phi=1$, and the resulting
lower bound. Source: accepted artefact \texttt{proof.md} (task B-C3-002), sections 0--6, gated VALID by two
referees (B-C3-008, B-C3-009); sha256\\ {\small\texttt{e86838d6c354e239e3201105542c8b83e77e6aa312dfc62488c6a18818c86c7b}}.
\end{itemize}

\setcounter{section}{-1}
\section{Notation}

\textbf{0.1.} We identify a partition $\lambda=(\lambda_1,\dots,\lambda_s)$ with its diagram
$C(\lambda)=\{(i,j)\in\N^2: 0\le j<\lambda_{i+1}\}$. Rows $i=0,\dots,s-1$ and columns $j$ are 0-based. A finite set
$S\subset\N^2$ is the diagram of a partition iff it is an \emph{order ideal}: $(i,j)\in S$ implies $(i-1,j)\in S$
when $i\ge1$ and $(i,j-1)\in S$ when $j\ge1$. Its rows are then initial segments, and their lengths are weakly
decreasing.

\textbf{0.2.} The \emph{diagonal index} of $(i,j)$ is $i+j$. Diagonal $d$ consists of the $d+1$ positions
$(i,d-i)$, $0\le i\le d$; we call $i$ the \emph{row} of the position. $D_m=\{(i,j):i+j\le m\}$ is the diagram of
$\delta_{m+1}$ ($D_{-1}=\emptyset$).

\textbf{0.3.} The energy is $E(S)=\sum_{(i,j)\in S}(i+j)$, and $E(\lambda)=E(C(\lambda))$.

\textbf{0.4.} Throughout, $n=T_{k-1}+r$ with $k\ge 2$ and $1\le r\le k$; this $k$ is the rank of $n$.

\section{Rotation lemma (R1)}

\textbf{1.1.} Define $R$ on $\N^2$ by $R(i,j)=(i+1,j-1)$ if $j\ge1$, and $R(i,0)=(0,i)$. It preserves $i+j$. On
diagonal $d$ it sends row $i$ to row $i+1$ for $i<d$ and row $d$ to row $0$, i.e.\ row $i\mapsto (i+1)\bmod (d+1)$.
So $R$ is a bijection of each diagonal, hence injective on $\N^2$ and energy-preserving: $E(R(S))=E(S)$.

\textbf{1.2.} Let $\lambda$ have $s$ parts and $C=C(\lambda)$. Row $0$ of $R(C)$ is
$\{(0,i):(i,0)\in C\}=\{(0,0),\dots,(0,s-1)\}$, of length $s$. For $0\le i\le s-1$, row $i+1$ of $R(C)$ is
$\{(i+1,j-1):1\le j<\lambda_{i+1}\}$, i.e.\ columns $0,\dots,\lambda_{i+1}-2$, of length $\lambda_{i+1}-1$. So every
row of $R(C)$ is an initial segment, and the row lengths in order are $\ell=(s,\lambda_1-1,\dots,\lambda_s-1)$.

\begin{lemma}[1.3]\label{lem:R1}
(i) If $\lambda_1\le s+1$, then $\ell$ is weakly decreasing, $R(C)$ is an order ideal, and $R(C)=C(B(\lambda))$.
So $E(B(\lambda))=E(\lambda)$. (ii) If $\lambda_1\ge s+2$, then $R(C)$ is not an order ideal and
$E(B(\lambda))\le E(\lambda)-1$.
\end{lemma}

\begin{proof}
(i) $s\ge\lambda_1-1\ge\lambda_2-1\ge\cdots\ge\lambda_s-1\ge0$. So the rows (initial segments) have weakly
decreasing lengths, which makes $R(C)$ an order ideal. The zero-length rows come last. By definition $B(\lambda)$
is the list of positive entries of $\ell$ sorted decreasingly, which here is $\ell$ with its trailing zeros
removed. So $C(B(\lambda))=R(C)$, and $E$ is preserved by 1.1.

(ii) Row 1 of $R(C)$ has length $\lambda_1-1>s$, so $(1,s)\in R(C)$ but $(0,s)\notin R(C)$: not an order ideal.
For an arrangement $\ell=(\ell_0,\ell_1,\dots)$ of row lengths, the set of left-justified rows has energy
$\sum_i\big(i\,\ell_i+\ell_i(\ell_i-1)/2\big)$. The second sum does not depend on the order, and zero rows
contribute $0$. Let $\ell'$ be $\ell$ with entries $0,1$ swapped. Then
$\sum_i i\ell'_i=\sum_i i\ell_i-(\lambda_1-1-s)\le\sum_i i\ell_i-1$. By the rearrangement inequality, the
decreasing arrangement $\ell^{\downarrow}$ of $\ell'$ satisfies $\sum_i i\ell^{\downarrow}_i\le\sum_i i\ell'_i$.
Now $C(B(\lambda))$ is the left-justified set with row lengths $\ell^{\downarrow}$ (zeros dropped at the end).
Hence $E(B(\lambda))\le E(R(C))-1=E(\lambda)-1$.
\end{proof}

\textbf{1.4. Corollary.} $E(B(\lambda))\le E(\lambda)$ always, and $E(B(\lambda))=E(\lambda)$ iff $R(C(\lambda))$
is an order ideal, in which case $C(B(\lambda))=R(C(\lambda))$. (Combine (i) and (ii); the two cases are
exhaustive.)

\section{Minimum energy and cyclicity (R2)}

\textbf{2.1.} For finite $S\subset\N^2$ with $|S|=n$, let $N_{\ge d}(S)=\#\{x\in S:\diag(x)\ge d\}$. Then
$E(S)=\sum_{x\in S}\diag(x)=\sum_{d\ge1}N_{\ge d}(S)$, since a cell on diagonal $e$ is counted once for each
$d=1,\dots,e$.

\textbf{2.2.} There are exactly $T_d$ positions with diagonal $\le d-1$. So $N_{\ge d}(S)\ge n-T_d$, and also
$N_{\ge d}\ge0$. Put $\nu_d=\max(0,n-T_d)$ and $g_d(S)=N_{\ge d}(S)-\nu_d\ge 0$. Since $T_d\le T_{k-1}<n$ for
$d\le k-1$ and $T_d\ge T_k\ge n$ for $d\ge k$, we get $\nu_d=n-T_d$ for $1\le d\le k-1$ and $\nu_d=0$ for $d\ge k$.

\textbf{2.3.} Set $E_{\min}(n)=\sum_{d\ge1}\nu_d$ and $\Phi(S)=E(S)-E_{\min}(n)=\sum_{d\ge1}g_d(S)$, a sum of
non-negative integers.

\textbf{2.4. Lemma.} $\Phi(S)=0$ iff $S=D_{k-2}\cup Q$ with $Q$ a set of $r$ positions on diagonal $k-1$.

\begin{proof}
$\Phi=0$ iff $g_d=0$ for all $d$. For $d\ge k$ this says $N_{\ge k}=0$: no cell has diagonal $\ge k$. For
$1\le d\le k-1$ it says $\#\{x:\diag(x)\le d-1\}=T_d$, i.e.\ all positions of diagonals $0,\dots,d-1$ are in $S$.
Taking $d=k-1$, $D_{k-2}\subseteq S$. The remaining $n-T_{k-1}=r$ cells lie on diagonal $k-1$. The converse is the
same computation read backwards.
\end{proof}

\begin{lemma}[2.5]\label{lem:R2}
$\lambda\vdash n$ is cyclic iff $\Phi(\lambda)=0$.
\end{lemma}

\begin{proof}
By Cell 1 (cited), the cyclic partitions are $\lambda(e)=(k-1+e_1,\dots,1+e_{k-1},e_k)$ with $e\in\{0,1\}^k$ and
$\sum e=r$. Row $i$ ($0\le i\le k-1$) of $\lambda(e)$ has length $k-1-i+e_{i+1}$. So
$C(\lambda(e))=D_{k-2}\cup\{(i,k-1-i):e_{i+1}=1\}$, which has $\Phi=0$ by 2.4. Conversely, suppose
$\Phi(\lambda)=0$. By 2.4, $C(\lambda)=D_{k-2}\cup Q$. Set $e_{i+1}=[(i,k-1-i)\in Q]$. Then row $i$ has length
$k-1-i+e_{i+1}$ for $0\le i\le k-1$ and there are no further rows. So $\lambda=\lambda(e)$ with $\sum e=r$, and
$\lambda$ is cyclic by Cell 1.
\end{proof}

\textbf{2.6. Consequence.} If $\Phi(\lambda)=1$, then $\lambda$ is not cyclic. Either $\Phi(B\lambda)=1$ and
$C(B\lambda)=R(C(\lambda))$, or $\Phi(B\lambda)=0$ and $B\lambda$ is cyclic. (By 1.4 the energy stays the same or
drops by $\ge 1$, and $\Phi\ge0$.)

\section{Partitions with $\Phi=1$ (R3)}

\textbf{3.1.} Let $\Phi(S)=1$ for an order ideal $S$ with $|S|=n$. Then there is exactly one $d^*$ with
$g_{d^*}=1$, and $g_d=0$ for $d\ne d^*$.

\textbf{3.2.} \emph{Case $d^*\ge k$.} Here $g_d=0$ for $d\le k-1$, so $D_{k-2}\subseteq S$ as in 2.4. For $d\ge k$,
$N_{\ge d}=g_d$. $N_{\ge d}$ is non-increasing in $d$ and $N_{\ge d^*}=1$, so if $d^*>k$ then $N_{\ge k}\ge1$,
contradicting $g_k=0$. Hence $d^*=k$, $N_{\ge k}=1$ and $N_{\ge k+1}=0$: exactly one cell lies off $D_{k-1}$, and
it is on diagonal $k$. Diagonal $k-1$ then holds $r-1$ cells. Call this \textbf{case B}.

\textbf{3.3.} \emph{Case $d^*\le k-1$.} Here $N_{\ge k}=0$. Suppose first $d^*\le k-2$. Then
$\#\{\diag\le d^*-1\}=T_{d^*}-1$ and $\#\{\diag\le d^*\}=T_{d^*+1}$ (because $g_{d^*+1}=0$). So diagonal $d^*$
would hold $d^*+2$ cells, but it has only $d^*+1$ positions. Hence $d^*=k-1$. Diagonals $0,\dots,k-3$ are full,
and diagonal $k-2$ has exactly one missing position $(i,k-2-i)$. The order-ideal property then excludes
$(i,k-1-i)$ (whose left neighbour is $(i,k-2-i)$) and $(i+1,k-2-i)$ (whose upper neighbour is $(i,k-2-i)$). So
diagonal $k-1$ holds at most $k-2$ cells. It holds $n-(T_{k-1}-1)=r+1$ cells, so $r\le k-3$. Call this
\textbf{case A}.

\begin{lemma}[3.4]\label{lem:R3}
If $r\ge k-2$, every $\lambda\vdash n$ with $\Phi(\lambda)=1$ is of case B. In particular this holds for
$n=T_k-1$, where $r=k-1$. (Immediate from 3.2 and 3.3.)
\end{lemma}

\textbf{3.5.} \emph{Parametrising case B.} Let $2\le r\le k-1$ and $h=k-r+1$, so $2\le h\le k-1$. A case-B diagram
is determined by the set $H\subseteq\{0,\dots,k-1\}$ of rows of the $h$ empty positions of diagonal $k-1$
(``holes'') and the row $c\in\{0,\dots,k\}$ of the cell on diagonal $k$:
\[
S(H,c)=D_{k-2}\cup\{(i,k-1-i):0\le i\le k-1,\ i\notin H\}\cup\{(c,k-c)\}.
\]
Call $(H,c)$ \emph{admissible} if ($c\ge1\Rightarrow c-1\notin H$) and ($c\le k-1\Rightarrow c\notin H$).

\textbf{Claim (R3$'$):} $S(H,c)$ is an order ideal iff $(H,c)$ is admissible.

\begin{proof}
A cell of $D_{k-2}$ has its up/left neighbours in $D_{k-3}\subseteq D_{k-2}$. A cell on diagonal $k-1$ has its
neighbours on diagonal $k-2$, which is full. The cell $(c,k-c)$ has upper neighbour $(c-1,k-c)$ (present iff
$c\ge1$ and $c-1\notin H$, required only if $c\ge 1$) and left neighbour $(c,k-1-c)$ (present iff $c\le k-1$ and
$c\notin H$, required only if $k-c\ge1$).
\end{proof}

The partition with diagram $S(H,c)$ is written $\Lambda(H,c)$. Its row $i$ ($0\le i\le k$) has length
$k-1-i+[i\notin H]+[i=c]$ (with row $k$ of length $[c=k]$).

\section{Exact dynamics in case B (R4)}

\textbf{4.1.} By 1.1, $R$ maps diagonal $d$ to itself by row $i\mapsto(i+1)\bmod(d+1)$. So
$R(S(H,c))=S(H+1,\,c')$, where $H+1=\{(a+1)\bmod k:a\in H\}$ and $c'=(c+1)\bmod(k+1)$. (It fixes $D_{k-2}$
setwise, permutes diagonal $k-1$ with period $k$ and diagonal $k$ with period $k+1$.)

\textbf{4.2.} Let $\lambda=\Lambda(H,c)$ with $(H,c)$ admissible. By 2.6 and 1.4: if $(H+1,c')$ is admissible, then
$R(C(\lambda))$ is an order ideal (3.5), so $B\lambda=\Lambda(H+1,c')$, again with $\Phi=1$. Otherwise
$R(C(\lambda))$ is not an order ideal, so the energy drops (1.4), so $B\lambda$ is cyclic (2.6). By induction on
$t$:
\[
d_B(\Lambda(H,c))=\tau(H,c):=\min\{t\ge1:\ (H+t,\ (c+t)\bmod(k+1))\text{ not admissible}\}.
\]
Here $H+t$ is taken mod $k$. For $t<\tau$ the partition $B^t\lambda=\Lambda(H+t,c+t)$ has $\Phi=1$ and is not
cyclic (2.5), and $B^{\tau}\lambda$ is cyclic. (Finiteness of $\tau$ follows from 4.4.)

\textbf{4.3.} \emph{Reduction to one hole.} For $\alpha\in[0,k-1]$ and $\gamma\in[0,k]$, the pair
$(\{\alpha\},\gamma)$ fails admissibility iff $\alpha=\gamma-1$ or $\alpha=\gamma$. The side conditions
$\gamma\ge1$ and $\gamma\le k-1$ are automatic in these two equalities, because $0\le\alpha\le k-1$. So
$(H,\gamma)$ is admissible iff every $(\{a\},\gamma)$, $a\in H$, is admissible. Hence
$\tau(H,c)=\min_{a\in H}\tau_a$, where $\tau_a$ is the first $t\ge1$ at which $w_t:=\alpha_t-\gamma_t\in\{-1,0\}$,
with $\alpha_t=(a+t)\bmod k$ and $\gamma_t=(c+t)\bmod(k+1)$.

\textbf{4.4. Lemma.} If $(\{a\},c)$ is admissible, then
$\tau_a=\min\{t\ge1: a+t\equiv0\pmod k,\ c+t\equiv1\pmod{k+1}\}=1-c+(k+1)\,m_a$, where
$m_a=(c-1-a)\bmod k\in[0,k-1]$.

\begin{proof}
Suppose $w_{t-1}\notin\{-1,0\}$. We have $w_{t-1}\in[-k,k-1]$. Four cases:
\begin{enumerate}\renewcommand{\labelenumi}{(\arabic{enumi})}
\item Neither $\alpha$ nor $\gamma$ wraps ($\alpha_{t-1}<k-1$, $\gamma_{t-1}<k$). Then $w_t=w_{t-1}\notin\{-1,0\}$.
\item Only $\alpha$ wraps ($\alpha_{t-1}=k-1$, $\gamma_{t-1}\le k-1$). Then $w_t=w_{t-1}-k$. Since $w_{t-1}\le k-1$,
this is $\le-1$, and it equals $-1$ iff $w_{t-1}=k-1$, i.e.\ $\gamma_{t-1}=0$, i.e.\ $(\alpha_t,\gamma_t)=(0,1)$.
\item Only $\gamma$ wraps ($\gamma_{t-1}=k$, $\alpha_{t-1}\le k-2$). Then $w_t=\alpha_{t-1}+1\ge1$.
\item Both wrap. Then $w_{t-1}=k-1-k=-1$, which is excluded.
\end{enumerate}
So if $t$ is the first failure time, $w_{t-1}$ is admissible (for $t=1$ by hypothesis), and only case (2) with
$(\alpha_t,\gamma_t)=(0,1)$ can fail. Conversely $(\alpha_t,\gamma_t)=(0,1)$ gives $w_t=-1$, a failure. Hence
$\tau_a$ is the least $t\ge1$ with $a+t\equiv0\pmod k$ and $c+t\equiv1\pmod{k+1}$.

Solve by the Chinese remainder theorem ($\gcd(k,k+1)=1$). Write $t=1-c+(k+1)m$. Since $k+1\equiv1\pmod k$, the
first congruence becomes $m\equiv c-1-a\pmod k$. So the solutions are $t=1-c+(k+1)(m_a+kq)$, $q\in\mathbb Z$. Put
$t_0=1-c+(k+1)m_a$. Then $t_0\le 1+(k+1)(k-1)=k^2$, so $t_0-k(k+1)\le -k<1$. And $t_0\ge1$: if $m_a\ge1$ then
$(k+1)m_a\ge k+1>c$. If $m_a=0$ then $a\equiv c-1\pmod k$; for $c\ge1$ this means $a=c-1$, which contradicts
admissibility, so $c=0$ and $t_0=1$. Hence $\tau_a=t_0$.
\end{proof}

\begin{lemma}[4.5]\label{lem:R4}
For admissible $(H,c)$:
\[
d_B(\Lambda(H,c))=1-c+(k+1)\min_{a\in H}\big((c-1-a)\bmod k\big).
\]
(From 4.2, 4.3 and 4.4.)
\end{lemma}

\begin{remark}
Exact check (not load-bearing): \texttt{code/check\_c3.py} compares this formula with brute-force $d_B$ for every
case-B partition with $2\le r\le k-1$, for $3\le k\le K$ (output lines ``(L)''). With $K=10$ all match; see
``How to verify'' below.
\end{remark}

\section{Maximum over case B (R5)}

\textbf{5.1.} The map $a\mapsto m_a=(c-1-a)\bmod k$ is a bijection of $\{0,\dots,k-1\}$, so the $m_a$ ($a\in H$)
are $h$ distinct values. Admissibility says: $c\ge1\Rightarrow m_a\ne0$ (since $a\ne c-1$), and
$c\le k-1\Rightarrow m_a\ne k-1$ (since $a\ne c$ iff $c-1-a\not\equiv -1$).

\textbf{5.2.} The minimum of $h$ distinct integers in an interval $[x,y]$ is at most $y-h+1$. Hence:
\begin{itemize}
\item $c=0$: $m_a\in[0,k-2]$, so $\min m_a\le k-1-h$ and $d_B\le 1+(k+1)(k-1-h)$.
\item $1\le c\le k-1$: $m_a\in[1,k-2]$, so $d_B\le 1-c+(k+1)(k-1-h)\le (k+1)(k-1-h)$.
\item $c=k$: $m_a\in[1,k-1]$, so $\min m_a\le k-h$ and $d_B\le 1-k+(k+1)(k-h)=(k+1)(k-1-h)+2$.
\end{itemize}

\begin{lemma}[5.3]\label{lem:R5}
Over all case-B partitions of $n=T_{k-1}+r$ ($2\le r\le k-1$), $\max d_B=(k+1)(r-2)+2$. It is attained only by
$c=k$ with $\{m_a\}=\{k-h,\dots,k-1\}$, i.e.\ $H=\{a=k-1-m\}=\{0,\dots,h-1\}$.
\end{lemma}

\begin{proof}[Justification (as in the accepted artefact)]
This $(H,c)$ is admissible: it needs $k-1\notin H$, i.e.\ $h\le k-1$, which holds. Here $h=k-r+1$ gives
$(k+1)(k-1-h)+2=(k+1)(r-2)+2$. The first two bullets of 5.2 are strictly smaller.
\end{proof}

\textbf{5.4.} For $n=T_k-1$ ($r=k-1$, $h=2$, $k\ge3$), the maximiser is $H=\{0,1\}$, $c=k$. Its row lengths (3.5)
are: row 0: $k-1$; row 1: $k-2$; row $i$ for $2\le i\le k-1$: $k-i$; row $k$: $1$. That is,
\[
\lambda^*_k=(k-1,\,k-2,\,k-2,\,k-3,\,\dots,\,2,\,1,\,1),
\]
with sum $(k-1)+(k-2)+T_{k-2}+1=T_k-1$. Its value is $(k+1)(k-3)+2=k^2-2k-1$. By 3.4 every $\Phi=1$ partition of
$T_k-1$ is of case B. So \textbf{$\lambda^*_k$ is the unique partition of $T_k-1$ with $\Phi=1$ and
$d_B=k^2-2k-1$.}

\section{Lower bound for (b) (R6)}

\begin{lemma}[6.1]\label{lem:R6}
For every $k\ge3$, $d_B(\lambda^*_k)=k^2-2k-1$ (5.4). So $D_B(T_k-1)\ge k^2-2k-1$.
\end{lemma}

\begin{proof}[Direct check with 4.5]
$m_0=k-1$ and $m_1=k-2$, so $d=1-k+(k+1)(k-2)=k^2-2k-1$. Examples: $k=4$: $(3,2,2,1,1)$, $d=7$. $k=5$:
$(4,3,3,2,1,1)$, $d=14$.
\end{proof}

\section*{Open parts}

The following are \textbf{OPEN}; no proof of them is part of this submission, and no conjectured answer is
presented as established.
\begin{itemize}
\item[(a)] The upper bound $d_B(\lambda)\le k^2-2k-1$ for every non-triangular $n$ of rank $k\ge4$ (every
$T_{k-1}<n<T_k$ and every $\lambda\vdash n$): OPEN.
\item[(b)] The upper half of (b), $D_B(T_k-1)\le k^2-2k-1$; so the equality $D_B(T_k-1)=k^2-2k-1$ is NOT
established: OPEN. Only the lower bound $D_B(T_k-1)\ge k^2-2k-1$, $k\ge3$ (Lemma~R6), is proved.
\item[(c)] The set of partitions of $T_k-1$ attaining $D_B(T_k-1)$: OPEN.
\end{itemize}

\paragraph{Numerical evidence (not proof).} The shipped exact checker \texttt{code/check\_c3.py}, run as
\texttt{python3 check\_c3.py 10}, enumerates every partition of every $n$ with $T_{k-1}<n<T_k$, $2\le k\le 10$, and
computes $d_B$ by exact cycle detection. Its output (reproduced below from the run recorded in
\texttt{verify/logs/check\_c3\_K10.txt}) reports $D_B(n)\le k^2-2k-1$ for all such $n$ with $4\le k\le 10$, and
$D_B(T_k-1)=k^2-2k-1$ for $4\le k\le 10$. This concerns exactly those finitely many $n$; it proves nothing for
$k\ge 11$ and is not part of the proof of any claim above.

\begin{center}
\begin{tabular}{c|ccccccccc}
$k$ & 2 & 3 & 4 & 5 & 6 & 7 & 8 & 9 & 10\\ \hline
$T_k-1$ & 2 & 5 & 9 & 14 & 20 & 27 & 35 & 44 & 54\\
$D_B(T_k-1)$ (computed) & 0 & 3 & 7 & 14 & 23 & 34 & 47 & 62 & 79\\
$k^2-2k-1$ & $-1$ & 2 & 7 & 14 & 23 & 34 & 47 & 62 & 79
\end{tabular}
\end{center}
For $k=2,3$ the computed values differ from $k^2-2k-1$. The checker's final line was \texttt{ALL OK}.

\section*{How to verify}

\begin{itemize}
\item Command (stdlib-only Python 3, exact integer arithmetic), from the directory \texttt{submission/code/}:\\
\texttt{/usr/bin/time -p python3 check\_c3.py 10}
\item Expected output: last line \texttt{ALL OK}.
\item Measured runtime: \texttt{real 26.93} s (\texttt{user 26.66}, \texttt{sys 0.14}), Python 3.14.0,\\ output \texttt{ALL OK}; well under the 10-minute limit.
\item The check is not load-bearing for Lemmas R1--R6; it checks Lemma~R4's formula (lines ``(L)'') and
$d_B(\lambda^*_k)=k^2-2k-1$ for $4\le k\le 10$ only.
\end{itemize}

\section*{Limitations}

\begin{itemize}
\item Only the lower half of (b) is established. (a), the upper half of (b) and (c) are OPEN.
\item Lemma~R2, and hence everything after it, relies on the Cell 1 classification of cyclic partitions, cited as
an earlier cell result and not re-proved here.
\item The accepted artefact also contains further material (its sections 7--9: an explicit family of partitions
with $d_B=k^2-2k-1$, partial upper-bound analysis, and maximiser computations). That material was not gated and is
deliberately not included here as a result.
\item The computation above is numerical evidence over a finite range only.
\end{itemize}

\end{document}
